% !PRINTED 415
\weberpage{415}
worin sich das Productzeichen $\Pi$ auf alle Theiler $\mu$ von $n$
bezieht.

Daraus erhalten wir nebenbei einen Beweis des zahlentheoretischen
Satzes
\weq{4}{n = \overset{\mu}{\Sigma}\varphi(\mu),}
wenn wir den Grad der Functionen auf der rechten und der
linken Seite von (3) einander gleich setzen.

Andererseits schliessen wir nach dem \gesperrt{Gauss'schen} Theorem
(\S.\,2), dass die Co\"efficienten der s\"ammtlichen $X_\mu$ \gesperrt{ganze Zahlen}
sind. Denn k\"amen darunter auch gebrochene Zahlen vor, so
k\"onnte das Product nicht lauter ganzzahlige Co\"efficienten enthalten, wie es doch nach (3) und (1) sein muss.

Alle Functionen $f_n(x)$ haben den Theiler $x-1$, und wenn
wir die Theilung ausf\"uhren, so ergiebt sich
\weq{5}{\frac{x^n-1}{x-1} = x^{n-1} + x^{n-2} + \cdots + x + 1.}
Hieraus schliessen wir, dass, wenn $r$ irgend eine primitive
oder nicht primitive $n^{\text{te}}$ Einheitswurzel ist, mit alleiniger Ausnahme von $r=1$, immer
\weq{6}{1 + r + r^2 + \cdots + r^{n-1} = 0,}
w\"ahrend f\"ur $r=1$ die Summe auf der linken Seite von (6)
offenbar den Werth $n$ hat.

Wenn $n$ eine Primzahl ist, so giebt es ausser 1 keine imprimitiven $n^{\text{ten}}$ Einheitswurzeln, und daher ist, wenn $n$ eine Primzahl ist, was wir dadurch andeuten wollen, dass wir $p$ daf\"ur
setzen,
\weq{7}{X_p = x^{p-1} + x^{p-2} + \cdots + x + 1.}
Ebenso einfach l\"asst sich $X_n$ bilden, wenn $n$ eine \gesperrt{Primzahlpotenz} ist, also
\[
n = p^{\pi},
\]
oder, wenn wir zur Abk\"urzung $p^{\pi-1} = p'$ setzen:
\[
n = p'p.
\]
Die $n^{\text{ten}}$ Einheitswurzeln bestehen in diesem Falle aus den
primitiven $n^{\text{ten}}$ Einheitswurzeln und aus den $p'^{\text{ten}}$ Einheitswurzeln
und es ist also:
\weq{8}{X_n = \frac{x^{pp'}-1}{x^{p'}-1} = x^{p'(p-1)} + x^{p'(p-2)} + \cdots + x^{p'} + 1}
